\section{Investigation Process}

Digital forensic investigations can be conceptually divided into a set of phases.\\ 
These phases do not represent a rigid or universal procedure, but rather a \textbf{conceptual framework} that helps investigators understand how to approach an investigation and what subtleties arise during each step.
\medskip

\noindent
Different standards propose slightly different models. Some frameworks describe four phases, others merge or rename certain steps.

\begin{quotation}
  \textbf{Author's npte:} \textit{ Porfessor Atzeny underlines that the goal is not to memorize a specific structure, but to understand the \textbf{underlying investigative logic}.}
\end{quotation}

\noindent
Well-known standards and guidelines include:

\begin{itemize}
\item \textbf{ACPO Guidelines} (United Kingdom)
\item \textbf{NIST Digital Forensics Framework}
\item \textbf{ISO standards} related to forensic processes
\end{itemize}

\noindent
Although these frameworks differ in terminology or number of steps, they describe broadly similar investigative processes.
\medskip

\noindent
The investigation process can be described using five phases:

\begin{enumerate}
\item Identification
\item Collection
\item Acquisition
\item Evaluation (or Examination)
\item Presentation
\end{enumerate}

\subsection{1. Identification}

The identification pphase consists of recognizing potential sources of digital evidence.\\
Typical sources include:

\begin{itemize}
\item storage devices (drives, removable media)
\item memory and running processes
\item mobile devices
\item network traffic
\item cloud storage and online services
\end{itemize}

\noindent
Investigatorsbegin forming preliminary hypotheses about relevance and also understanding the \textbf{volatility of information sources}. \\
Some evidence may disappear or change quickly, which means investigators must later prioritize acquisition accordingly.  

\medskip
\noindent
Whenever possible, investigators should gather contextual information \textbf{before interacting with the scene}, since direct interaction may alter digital evidence.
\\
Possible preliminary sources include:

\begin{itemize}
\item surveillance videos or photographs
\item social media content
\item organizational records or logs
\end{itemize}

\noindent
Another important factor is the evaluation of \textbf{technical knowledge of the suspect}.\\
If the suspect has limited technical expertise, the likelihood of sophisticated counter-forensic mechanisms may be lower.

\subsubsection{Scene Identification}

Once investigators reach the scene, the identification process continues with the recognition of actual devices and artifacts.\\
During this phase investigators may build a preliminary \textbf{graph of entities}, connecting devices, users, and network components that could be involved in the incident.

\subsubsection{Volatility and Prioritization}

Understanding volatility is essential because some evidence may disappear during normal system activity.
\\
Highly volatile sources should generally be acquired first.
\\
For instance:

\begin{itemize}
\item a powered-off device should normally remain powered off until acquisition
\item volatile memory may require specialized techniques to preserve its contents
\end{itemize}

\noindent
This prioritization begins during the identification phase and guides the subsequent collection process.

\subsubsection*{Example Scenario: Insider Data Exfiltration}

Consider an investigation inside a company where an employee is suspected of exfiltrating sensitive documents.\\
Potential evidence sources may include:

\begin{itemize}
\item employee workstations
\item USB storage devices
\item printers and shared devices
\item temporary files and browser caches
\item cloud storage services
\item internal file sharing systems
\end{itemize}

\noindent
Investigators may begin by collecting publicly available information about employees or suspects.\\  
Tools commonly used for this type of intelligence gathering include:

\begin{itemize}
\item \textbf{SpiderFoot}
\item \textbf{Maltego}
\item \textbf{Shodan}
\end{itemize}

\noindent
These tools help correlate information.
\medskip

\noindent
In managed environments investigators may also collect system information through administrative access.
\\
Typical diagnostic commands can reveal useful information about system activity, including:

\begin{itemize}
\item connected USB devices
\item mounted storage devices
\item open files and processes
\item active network connections
\end{itemize}

\noindent
For istance:

\begin{itemize}
\item \texttt{lsusb}, \texttt{mount}, \texttt{lsof}, \texttt{netstat} on Unix-like systems
\item administrative tools or PowerShell commands on Windows systems
\end{itemize}


\subsection{ 2. Collection Phase}

Once potential evidence sources have been identified, the next step is the \textbf{collection phase}.\\ 
The goal is to secure and isolate the identified evidence sources while minimizing any modification to their content.
\medskip
\noindent
A fundamental rule is to \textbf{avoid unnecessary interaction with the system}.

\subsubsection{Isolation of Evidence}

A central objective of the collection phase is the \textbf{isolation of evidence sources} to prevent further modification or destruction.\\
Isolation may occur at two levels:

\begin{itemize}
\item \textbf{Physical isolation}: preventing the suspect from accessing the device.
\item \textbf{Logical isolation}: preventing network communication with the device.
\end{itemize}

\noindent
Logical isolation may involve modifying network configurations such as firewall rules, network topology, or gateway settings to block communication with the device.

\subsubsection{Preservation of Evidence}

Investigators must also protect evidence from environmental or physical damage.\\
For examples:

\begin{itemize}
\item protecting magnetic storage devices from strong magnetic fields
\item handling fragile hardware carefully
\item preserving volatile memory when possible
\end{itemize}

In advanced scenarios, volatile memory may even be preserved using specialized techniques such as cryogenic preservation to delay data loss.

\subsubsection{Documentation and Evidence Tracking}

During the collection phase, all actions must be carefully documented as part of the chain of custody.\\

This documentation typically includes:

\begin{itemize}
\item who performed the action
\item when the action occurred
\item what operation was performed
\item which device was involved
\end{itemize}

\noindent
Devices should also be uniquely identified using information such as:

\begin{itemize}
\item model number
\item serial number
\item physical description
\end{itemize}

\subsubsection*{Forensic Equipment}

Forensic investigators typically use specialized equipment during the collection phase.\\
Examples include:

\begin{itemize}
\item hardware write blockers
\item forensic imaging devices
\item storage adapters for disk acquisition
\item signal jammers to prevent wireless communication
\end{itemize}
